\subsection{Purpose}

This appendix fixes a minimal formalization of the subgraph contract and the \texttt{ComputeSubgraphContract} algorithm used for decomposition, aggregation, and local verification of the structural realization.

\subsection{Preconditions}

Let the structural realization be given as a typed directed multigraph


$$G = (V, E, \mathrm{src}, \mathrm{dst}, \tau_V, \tau_E),$$

where $\tau_V(v) \in \{\texttt{function}, \texttt{table}\}$ and $\tau_E(e) \in \{\texttt{consumes}, \texttt{produces}, \texttt{reads}, \texttt{writes}, \texttt{triggers}\}$.

Each function $f$ has a contract


$$C(f) = \langle C_{\mathrm{in}}, C_{\mathrm{out}}, R, W, T_{\mathrm{in}}, T_{\mathrm{out}} \rangle,$$

where all components are fully derived from incident edges of the corresponding type.

\subsection{Subgraph contract}

For a chosen subgraph $S$, the set of its vertices is denoted by $V(S)$ and decomposed into two parts:


$$V(S) = V_{\mathrm{func}}(S) \cup V_{\mathrm{table}}(S),$$

where $V_{\mathrm{func}}(S)$ is the set of vertices of type \texttt{function} belonging to $S$, and $V_{\mathrm{table}}(S)$ is the set of vertices of type \texttt{table} belonging to $S$.


$$E(S) = \{ e \in E(G) \mid \mathrm{src}(e) \in V(S) \lor \mathrm{dst}(e) \in V(S) \}$$

is the set of edges of graph $G$ for which at least one endpoint belongs to $V(S)$.

For the subgraph $S$, the contract


$$C(S) = \langle C_{\mathrm{in}}(S), C_{\mathrm{out}}(S), R(S), W(S), T_{\mathrm{in}}(S), T_{\mathrm{out}}(S) \rangle$$

is defined. It describes only the external boundary of the subgraph: which tables and functions connect $S$ to the rest of the graph.

\subsection{\texttt{ComputeSubgraphContract} algorithm}

\begin{lstlisting}
function ComputeSubgraphContract(S):
    V_func(S)  = {v in V(S) : t_V(v) = function}
    V_table(S) = {v in V(S) : t_V(v) = table}

    internal_tables = {t in V_table(S) :
        for all e in E where (src(e)=t or dst(e)=t) :
            (src(e) in V(S) and dst(e) in V(S))}

    interface_tables = V_table(S) \ internal_tables

    Cin(S)  = {t in interface_tables : exists e in E(S),
               t_E(e) = consumes,
               src(e) = t, dst(e) in V_func(S)}

    Cout(S) = {t in interface_tables : exists e in E(S),
               t_E(e) = produces,
               src(e) in V_func(S), dst(e) = t}

    R(S)    = {t in interface_tables : exists e in E(S),
               t_E(e) = reads,
               src(e) = t, dst(e) in V_func(S)}

    W(S)    = {t in interface_tables : exists e in E(S),
               t_E(e) = writes,
               src(e) in V_func(S), dst(e) = t}

    Tin(S)  = {f not in V(S) : exists e in E,
               t_E(e) = triggers,
               src(e) = f, dst(e) in V_func(S)}

    Tout(S) = {f not in V(S) : exists e in E,
               t_E(e) = triggers,
               src(e) in V_func(S), dst(e) = f}

    return <Cin(S), Cout(S), R(S), W(S), Tin(S), Tout(S)>
\end{lstlisting}

\subsection{Note on interface tables}

An interface table may belong to several components of the subgraph contract at once. For example, a table may be consumed by one boundary function and produced by another. This is a normal case: the subgraph contract fixes external relations rather than imposing a single role on a table.

\subsection{Usage}

\subsubsection{Decomposition}

Under decomposition, the node $f$ is replaced with a subgraph $S$ such that the following holds:


$$C(f) = C(S).$$

Intermediate internal tables are allowed as long as they do not reach the boundary of the subgraph.

\subsubsection{Aggregation}

Under aggregation, the subgraph $S$ is collapsed into a single node $f_S$ with the contract


$$C(f_S) = C(S).$$

This makes \texttt{ComputeSubgraphContract} the minimal formal mechanism that supports contract-preserving transformations between different levels of detail of the same structural realization.
